% !PRINTED 546
\weberpage{546}

\begin{center}
§. 165.

\gesperrt{Solution of the Cyclic Equations.}
\end{center}

The \gesperrt{Lagrange} resolvents, by application of the theorems now proved, lead to the solution of the cyclic equations — more precisely, to the reduction to \gesperrt{pure equations}.

We now understand by the $\alpha$ no longer arbitrary variables, but the roots of a cyclic equation, so that the cyclic functions of the $\alpha$ are to be regarded as known quantities.

By the theorem of \S.\,164, 2., the coefficients of $(\varepsilon,\alpha)^m$ are cyclic functions of the $\alpha$.

If we understand by $a_0, a_1\dots a_{m-1}$ quantities in $\Omega$ and set
\weq{1}{\psi_\lambda = a_0 + a_1\varepsilon^\lambda + a_2\varepsilon^{2\lambda} + \cdots + a_{m-1}\varepsilon^{(m-1)\lambda},}
then from this theorem it follows that
\weq{2}{(\varepsilon^\lambda,\alpha) = \sqrt[m]{\psi_\lambda}.}

If $\varepsilon$ denotes therein a primitive $m$th root of unity, then all the resolvents are contained in the form $(\varepsilon^\lambda,\alpha)$.

If we observe further that $(1,\alpha) = a$, as the sum of the roots, belongs to the known quantities, then by \S.\,164 (3) we have
\weq{3}{m\alpha = a + \sqrt[m]{\psi_1} + \sqrt[m]{\psi_2} + \cdots + \sqrt[m]{\psi_{m-1}},}
and with this $\alpha$ is expressed by radicals of degree $m$, which under the root signs contain, besides the quantities that occur in $\Omega$ from the outset, also $m$th roots of unity.

Each of these radicals has, considered by itself, $m$ different values, which differ from one another by $m$th roots of unity as factors. If we give each $m$th root all its values, we obtain from (3) many different values of $\alpha$, among which, by \S.\,164, (4), all the roots $\alpha, \alpha_1, \alpha_2\dots\alpha_{m-1}$ occur. But the number of the expressions thus derived from (3) is much greater, and it remains to distinguish those which are to be retained from those which are to be discarded. The following way leads to this most simply.
